**Conventions.** All graphs are finite, simple, undirected.
For graphs (A,B), the (two–color) Ramsey number (R(A,B)) is the least integer (N) such that for every red/blue coloring of the edges of the complete graph (K_N), there is either a red copy of (A) or a blue copy of (B). Equivalently: for every graph (F) on (N) vertices, either (F) contains (A) or (\overline{F}) contains (B).

**Asymptotic notation.** For functions of an integer parameter (n), (f(n)\ll g(n)) means: (\exists C>0) (allowed to depend on the fixed graph (G) but **not** on (n) or on the choice of the varying graph(s)) such that (f(n)\le C g(n)) for all relevant (n).

**Given.** A fixed graph (G) satisfying:

1. (**Tree condition**) (\exists C_T>0) such that for every (n\ge1) and every tree (T) on (n) vertices,
   [
   R(G,T)\le C_T, n.
   ]
2. (**Clique condition**) (\exists C_K>0) such that for every (n\ge1),
   [
   R(G,K_n)\le C_K, n^2.
   ]

**Question.** Does it follow that:

3. (**Size-linear conclusion**) (\exists C>0) such that for every graph (H) with no isolated vertices and with (m=e(H)\ge1),
   [
   R(G,H)\le C, m , ?
   ]

This is exactly the property that (G) is **Ramsey size-linear** (also phrased as (R(G,F)=O(e(F))) for every graph (F) with no isolated vertices).
